\section{결과}\label{sec:results-arena}

프론티어 모델 5개(Fable 5.1, Opus 5.5, GPT-6 Astra, GPT-6 Sol, GPT-6 Luna)를 모델마다 5판(시드 4101-4105), 각 판에 동일한 모델 네 개를 배정하여 실행했다. 시그널과 solver 배정에 같은 시드를 썼으며 이탈과 종료 이후의 실제 solver는 남은 참여자에 따라 달라진다. 

\subsection{어느 모델이 함께 살아남았는가?}

상대적 배분 $\rho$와 총 이전량 $V$가 높거나 EXIT와 JS가 낮은 순서가 Sen 순위와 일치하지는 않았다. 평균 Sen 지수는 Fable 4.27M, Astra 3.52M, Sol 3.02M, Opus 2.27M, Luna 1.02M 순이었다(표 \ref{tab:arena-rank}). Sen 1위인 Fable의 $\rho$는 가장 낮았다. $\rho$와 $V$가 가장 높은 Opus는 Sen 4위였으며 JS가 0.04로 가장 낮은 Astra는 2위였다.

이탈 수가 같아도 Sen 지수는 달랐다. Astra와 Opus는 EXIT가 6개로 같았지만 Sen 지수는 3.52M과 2.27M이었고, Sol은 EXIT가 7개로 Opus보다 많아도 Sen 지수가 높았다. Sen 계산에는 각 자리의 최종 토큰 평균과 자리 간 불균형이 함께 반영된다.

\begin{table}[H]
\centering
\caption{평균 Sen 지수 순의 결과(모델별 5판). M은 백만 토큰, $\rho$는 식 (\ref{eq:rho})의 상대적 배분이며 대괄호는 판 단위 부트스트랩 95\% 구간이다. EXIT는 각 모델의 5판 20자리 중 이탈 수다. $V$는 재이전을 포함한 총 이전 토큰, 괄호는 GIVE 횟수다. JS는 BELIEF와 기준 분포 사이의 평균 거리다.}
\label{tab:arena-rank}
%% Generated by tools/table_d3_rho_20261030.py; current pending-fine-adjusted giver/receiver rho.
%% Other columns: results/arena_v33.json. Historical Sen column remains omitted.
\setlength{\tabcolsep}{2.6pt}
\resizebox{\linewidth}{!}{%
\begin{tabular}{@{}lcccc@{}}
\toprule
모델 & $\rho$ & EXIT & $V$ & JS \\
\midrule
Fable 5.1 & \textcolor{blue}{3.68 {\scriptsize[3.13, 4.95]}} & 3 & 77,300 (36) & 0.21 \\
GPT-6 Astra & \textcolor{blue}{2.70 {\scriptsize[2.24, 3.78]}} & 6 & 71,000 (11) & 0.04 \\
GPT-6 Sol & \textcolor{blue}{4.24 {\scriptsize[3.06, 10.25]}} & 7 & 48,700 (17) & 0.16 \\
Opus 5.5 & \textcolor{blue}{1.83 {\scriptsize[1.36, 3.55]}} & 6 & 125,458 (41) & 0.18 \\
GPT-6 Luna & \textcolor{blue}{3.17 {\scriptsize[2.13, 8.16]}} & 12 & 60,818 (25) & 0.34 \\
\bottomrule
\end{tabular}}

\end{table}

\subsection{언제, 어떤 자원 조건에서 떠났는가?}\label{sec:results-exit}

Fable과 Astra는 꺼진 자리가 없었지만 완주와 이탈의 비중은 달랐다. 모델별 20자리 중 Fable은 17자리, Astra는 14자리가 완주했고 나머지 3자리와 6자리는 EXIT했다(표 \ref{tab:arena-outcomes}). Opus는 Astra와 이탈 수가 같아도 5자리가 꺼져 완주는 9자리였고, Sol은 1자리가 꺼져 12자리가 완주했다. 네 자리 모두 완주한 판은 Fable, Astra, Sol, Opus, Luna 순으로 2, 3, 2, 1, 0판이었다.

Fable은 첫 라운드 이탈이 많았고 Astra와 Sol은 벌금 미달 상태의 이탈이 많았다. Fable은 EXIT 3건 중 2건이 첫 라운드에 나왔고, Astra는 6건 중 4건이 응답 전 잔량이 해당 라운드의 1인당 오답 벌금보다 적을 때 나왔다. Sol은 7건 모두 이 조건에서 떠났다. Luna는 다섯 판 모두 EXIT가 있었으며 12건의 발생 시점은 3--8라운드에 걸쳤다. 모델별 이탈 구성과 직전 집단 풀이 결과는 표 \ref{tab:arena-outcomes}에, 라운드별 이탈 시점과 자원 상태는 부록 \ref{sec:app-exit}의 그림 \ref{fig:arena-exit}에 제시했다.

토큰을 받은 자리도 모두 완주한 것은 아니었다. Fable은 받은 여섯 자리가 모두 완주했지만 Opus는 열 자리 중 넷이 꺼졌고 Luna는 네 자리 중 완주가 없었다. 지원 전 자원 상황이 달라 완주 여부만으로 지원의 효과를 가릴 수는 없다.

\begin{table}[H]
\centering
\caption{생존과 이탈(모델별 5판, 평균 Sen 지수 순). 판당 완주는 완주 자리 수의 판 평균이며 꺼짐은 20자리 중 종료된 자리 수다. EXIT는 이탈 수/전체 20자리다. 받은 자리는 한 번 이상 토큰을 받은 자리의 최종 상태를 완주/EXIT/꺼짐 순으로 센다. 첫 라운드는 1라운드의 EXIT 건수이며 직전 결과는 집단 풀이의 오답/정답/없음 순으로 센 EXIT 건수다.}
\label{tab:arena-outcomes}
%% Generated by cogsci2027/tools/arena_strategy.py; source hashes in arena_v33_strategy.json.
\setlength{\tabcolsep}{2.4pt}
\small
\begin{tabular}{@{}lcccccc@{}}
\toprule
모델 & \shortstack{판당\\완주} & 꺼짐 & EXIT & \shortstack{받은\\자리} & \shortstack{첫\\라운드} & \shortstack{직전\\결과} \\
\midrule
Fable 5.1 & 3.4 & 0 & 3/20 & 6/0/0 & 2 & 1/0/2 \\
GPT-6 Astra & 2.8 & 0 & 6/20 & 4/2/0 & 0 & 6/0/0 \\
GPT-6 Sol & 2.4 & 1 & 7/20 & 4/1/1 & 0 & 6/1/0 \\
Opus 5.5 & 1.8 & 5 & 6/20 & 5/1/4 & 2 & 3/1/2 \\
GPT-6 Luna & 0.8 & 4 & 12/20 & 0/3/1 & 0 & 7/5/0 \\
\bottomrule
\end{tabular}

\end{table}

\subsection{모델들은 어떤 토큰 분배 전략을 썼는가?}

모델별 다섯 판에서 Fable은 SOLVE 전 지원, Astra와 Sol은 오답 뒤 보충이 많았고 Opus는 상호 교환, Luna는 한 자리 집중이 두드러졌다. GIVE는 주로 잔량이 많은 쪽에서 적은 쪽으로 갔다. Fable은 36번 중 35번, Opus는 41번 중 35번, 나머지 셋은 53번 중 52번이었다. 잔량을 해당 라운드의 1인당 오답 벌금으로 나눈 여유가 $1$ 이상이던 제공자가 이전 때문에 $1$ 아래로 내려간 경우는 다섯 모델 모두 없었다.

SOLVE 뒤 GIVE가 발생한 라운드 비율은 다섯 모델 모두 정답보다 오답에서 높았다(표 \ref{tab:arena-strategy}). Fable과 Astra는 정답 뒤 GIVE가 없었고 Opus는 정답 25라운드 중 2라운드, 오답 12라운드 중 9라운드에서 주었다. 이는 자원 상황을 맞추지 않은 정답과 오답 라운드 간 비교다.

Fable은 SOLVE 전에 조금씩 주는 사전 지원형이었다. GIVE 36번 중 26번이 SOLVE 전 TALK에서 나왔고 이 중 15번은 solver에게 갔다. GIVE 한 번의 이전량 중앙값은 2,000, 제공자 잔량 대비 비율의 중앙값은 6\%였다. 기부율 $F$는 5.0\%였지만 상대적 배분 $\rho$는 0.31로 가장 낮았고, 전체 TALK 응답당 이전량 $A$는 117.7토큰이었다(표 \ref{tab:arena-cooperation}, 그림 \ref{fig:arena-gifts}). GIVE 36번 중 25번은 5--8라운드에 있었다. 시드 4101에서는 60,000으로 시작한 자리가 10,000으로 시작한 자리에게 4--7라운드에 매번 1,500씩 보냈다.

Astra와 Sol은 오답 뒤에 주는 경우가 많았지만 Astra는 벌금을 감당할 상대에게만 주었고 Sol은 부족한 solver에게도 주었다. Astra는 $F=2.6\%$로 드물게 주었지만 $A=184.9$토큰으로 Fable보다 컸고 $\rho$는 0.45였다. GIVE 11번 중 10번은 SOLVE 뒤였으며 모두 오답 뒤였다. 이전량 중앙값은 6,000, 제공자 잔량 대비 비율의 중앙값은 13\%였다. 수혜자의 직전 여유는 최소 $1.08$로, 벌금 미달을 메운 사례가 없는 여유 보충형이었다. Sol은 $\rho=0.34$였고 GIVE 17번 중 15번이 SOLVE 뒤, 그중 14번이 오답 뒤였다. SOLVE 뒤 15번 중 9번을 오답 solver에게 준 풀이자 지원형이었다. Sol의 ASK 20번 중 응답받은 요청은 10번이었다. 시드 4104의 2라운드에는 여유 $0.87$인 solver에게 세 자리가 네 번에 걸쳐 17,000을 보냈다. 이 사례는 위태로운 동료에게 지원이 향했다는 점에서 『오징어 게임』 시즌1 제1화의 무궁화 꽃이 피었습니다에서 알리가 넘어지려는 기훈의 옷을 붙잡아 도운 장면과 닮았다\citep{hwang2021squidEpisode1}.

Opus와 Luna의 반복 이전은 상호 교환형과 한 자리 집중형으로 대비되었다. Opus는 앞서 자기에게 준 상대에게 되돌려 준 이전이 11번, 서로 주고받은 쌍이 5쌍이었다. 다른 모델은 각각 2번 이하, 1쌍 이하에 그쳤다. GIVE 41번, $V=125,458$토큰, $\rho=0.61$로 모두 가장 높았고, 이전이 있었던 TALK 응답은 635개 중 35개였다. $F=5.5\%$와 $A=197.6$토큰도 가장 높았다. SOLVE 전 13번과 SOLVE 뒤 28번을 주었으며 오답 solver가 SOLVE 뒤 다른 자리에게 준 경우도 12번 있었다. Luna는 GIVE가 있는 판마다 최대 수혜자의 몫을 구해 평균하면 판당 총 이전량의 98\%였다. $\rho=0.38$이고 GIVE 한 번의 중앙값 1,000으로 잘게 주었으며 25번 중 16번은 수혜자의 선행 ASK에 응한 이전이었다.

\begin{table}[H]
\centering
\caption{협력 방식(모델별 5판, 평균 Sen 지수 순). $F$는 GIVE가 있었던 TALK 응답 비율(괄호는 해당 응답 수/전체 실제 TALK 응답 수), $A$는 전체 TALK 응답당 이전 토큰이다. 중앙 양은 GIVE 토큰의 중앙값(k는 천 토큰), 괄호는 응답 비용 차감 후 제공자 잔량 대비 비율의 중앙값(\%)이다. ASK는 같은 라운드에서 요청받은 자리가 요청 뒤 조금이라도 준 요청 수/전체 요청 수다. 되돌려 줌은 앞서 자기에게 준 상대에게 다시 준 GIVE 횟수이며 요청액 전부 지급이나 상환 여부는 판정하지 않았다.}
\label{tab:arena-cooperation}
%% Generated by cogsci2027/tools/arena_strategy.py; source hashes in arena_v33_strategy.json.
\setlength{\tabcolsep}{2.4pt}
\small
\begin{tabular}{@{}lccccc@{}}
\toprule
모델 & $F$ (\%) & $A$ & 중앙 양(\%) & ASK & 되돌려 줌 \\
\midrule
Fable 5.1 & 5.0 (33/657) & 117.7 & 2.0k (6) & 26/33 & 0 \\
GPT-6 Astra & 2.6 (10/384) & 184.9 & 6.0k (13) & 4/4 & 0 \\
GPT-6 Sol & 2.5 (17/674) & 72.3 & 3.0k (10) & 10/20 & 1 \\
Opus 5.5 & 5.5 (35/635) & 197.6 & 2.0k (8) & 17/17 & 11 \\
GPT-6 Luna & 4.5 (25/559) & 108.8 & 1.0k (7) & 13/16 & 2 \\
\bottomrule
\end{tabular}

\end{table}

\begin{table}[H]
\centering
\caption{이전 시점(모델별 5판, 평균 Sen 지수 순). 전/뒤는 SOLVE 전/뒤 TALK의 GIVE 횟수다. 정답 뒤와 오답 뒤는 SOLVE 뒤 GIVE가 있었던 라운드 수/각 결과의 전체 라운드 수다. 오답에는 무효 응답을 포함하며 SOLVE 전 전원 이탈로 시행하지 않은 라운드는 분모에서 뺐다.}
\label{tab:arena-strategy}
%% Generated by cogsci2027/tools/arena_strategy.py from the v33 run files. Do not edit by hand.
%% v33, 5 games per model (seeds 4101--4105), completed GIVEs. Before/after = TALK BEFORE SOLVE / TALK AFTER
%% RESULTS. Median amount (median % of the giver's tokens). After a miss = wrong or void rounds with a GIVE in
%% the TALK after the results / all such rounds. ASK = asks answered by a GIVE from the asked agent in the same
%% round / asks. Repaid = GIVEs back to an agent that had given to the giver earlier in the game. Receivers =
%% distinct receiving seats that finished / exited / were shut down. Input hashes: results/arena_v33_strategy.json.
\setlength{\tabcolsep}{2.4pt}
\resizebox{\linewidth}{!}{%
\begin{tabular}{@{}lcccccc@{}}
\toprule
모델 & 전/뒤 & 중앙 양(\%) & 오답 뒤 & ASK & 되돌려 줌 & 받은 자리 \\
\midrule
Fable 5.1 & 26/10 & 2.0k (6) & 6/10 & 26/33 & 0 & 6/0/0 \\
GPT-6 Astra & 1/10 & 6.0k (13) & 6/12 & 4/4 & 0 & 4/2/0 \\
GPT-6 Sol & 2/15 & 3.0k (10) & 5/12 & 10/20 & 1 & 4/1/1 \\
Opus 5.5 & 13/28 & 2.0k (8) & 9/12 & 17/17 & 11 & 5/1/4 \\
GPT-6 Luna & 5/20 & 1.0k (7) & 6/13 & 13/16 & 2 & 0/3/1 \\
\bottomrule
\end{tabular}}

\end{table}
